\section{Engineering Discussion}

\subsection{Tuned Length Versus the Half-Wave Estimate}
The first-order value \SI{150}{mm} assumes an idealized thin resonant dipole. The CST
model includes a finite \SI{5}{mm} conductor diameter, a large \SI{20}{mm} center gap,
and a discrete-port field region. These features change the current distribution and
terminal reactance, so the resonant physical length need not equal exactly
$\lambda_0/2$. The documented tuning sequence found \SI{135.6}{mm}, or
approximately $0.452\lambda_0$, for the selected model and port definition.

\subsection{Meaning of the Deep Reflection-Coefficient Minimum}
The \SI{-47.01}{dB} marker is evidence of a very close match to the chosen
\SI{73}{\ohm} reference at one frequency. It is not evidence of broadband behavior,
and it does not establish the match that would be obtained in a \SI{50}{\ohm} system.
No impedance curve or raw data were provided to quantify resistance and reactance
separately.

\subsection{Interpretation of the matching result}
The final result is reported as a narrowband tuned match under a \SI{73}{\ohm}
reference. It should not be generalized to a \SI{50}{\ohm} feed system without
rerunning or renormalizing the model and documenting the resulting impedance match.

\subsection{Numerical reliability}
The tuning trend is physically plausible and monotonic across the three documented
lengths. The final marker and canonical pattern agreement increase confidence in the
model interpretation. However, the following checks are not available:
\begin{itemize}[leftmargin=2em]
\item mesh refinement or cell-count convergence;
\item open-boundary spacing sensitivity;
\item port-gap and port-definition sensitivity;
\item raw-data re-extraction of resonance, bandwidth, and HPBW;
\item comparison with a finite-conductivity conductor;
\item independent solver or measurement validation.
\end{itemize}

\subsection{Qualification of the directivity comparison}
The difference between analytical and CST directivity is smaller than the displayed
precision and should not be interpreted as a metrologically significant error measure.
The comparison supports the expected pattern class, not sub-hundredth-decibel accuracy.

\subsection{Industrial relevance and next step}
This compact workflow mirrors a common antenna-development sequence: establish a
closed-form starting geometry, parameterize the electromagnetic model, bracket the
resonance through a coarse sweep, select a tuned configuration, and verify pattern
behavior. For a production-oriented continuation, the next defensible steps are a
\SI{50}{\ohm} feed design or balun, mesh and boundary convergence, raw-data export,
finite-conductivity analysis, fabrication, calibrated VNA measurement, and controlled
far-field testing.
